\subsection{Primitive character vanishing at an arbitrary odd conductor}

\begin{claim}{A.C2.19}{Proved}
Let $K$ be finitely generated over $\F_2$, and let $C/K$
be smooth proper geometrically connected of genus two.
Suppose that, geometrically, a line bundle $L$ of exact
odd order $n>1$ has the following property: every
multiplication-by-two orbit on $(\Z/n)^\times$ meets
\[
                         \{i:L^i\in\Theta\}.
\]
Then the geometrically prime-to-two point-section map
of $C$ is injective. The assertion holds over every
finite extension, for forms satisfying this geometric
hypothesis, and for hyperbolic open subcurves.
There is no restriction on the number of prime
divisors of $n$, and $C$ need not be ordinary.

Equivalently, it suffices that the primitive
conductor-$n$ factor of the associated cyclic
\'etale cover have two-rank zero. All characters of
smaller conductor may survive.
\end{claim}
\begin{proof}
We first state and prove the field estimate used below.
If a subfield $M$ of a curve function field $F/k$ over
an algebraically closed field is generated by finitely
many functions each of degree $d$, and $e=[F:M]$,
then
\[
                    e\mid d,\qquad
                    g(M)\le(d/e-1)^2.                    \tag{AP.1}
\]
This is the controlled-field assertion of
\citet[Lemma 2.2 and Corollary 2.3, author p.~3]
{A_PrimPoonen2007}. For one generator it is immediate.
Suppose at an induction stage that $[F:M]=d/s$
and $g(M)\le(s-1)^2$. Adjoin another degree-$d$
function $h$, and put $a=[M(h):M]$. Then
\[
       [F:M(h)]=d/(as),\qquad [M(h):k(h)]=as.
\]
Castelnuovo--Severi gives
\[
 \begin{aligned}
 g(M(h))
 &\le a(s-1)^2+(a-1)(as-1)\\
 &=(as-1)^2-as(a-1)(s-1)\le(as-1)^2.
 \end{aligned}
\]
This proves the induction. In particular, the estimate
is valid for a compositum with many generators.

If $n$ is a prime power, A.C2.8 already proves the
claim. Suppose therefore that $r=\omega(n)\ge2$,
and put $d=2^r$. We have
\[
                            n>(d-1)^2.                   \tag{AP.2}
\]
For $r=2$, use $n\ge15>9$. For $r\ge3$,
$n\ge3\cdot5^{r-1}>4^r>(2^r-1)^2$:
the middle inequality holds at $r=3$ and its ratio
increases by a factor $5/4$ at each step.

Let $D\to C$ be the connected cyclic \'etale cover
associated with $L$, and let $G=\langle\sigma\rangle$
be its group. Its genus is $n+1$. The orbit hypothesis
and A.C2.5 kill all $\varphi(n)$ primitive characters,
so
\[
 f(D)\le f(C)+n-1-\varphi(n)
       \le n+1-\varphi(n)<n-1,                            \tag{AP.3}
\]
since $\varphi(n)\ge8$.
Define
\[
 a_\ell=\sigma^{n/\ell}\quad(\ell\mid n\text{ prime}),
 \qquad A=\prod_{\ell\mid n}(1-a_\ell).
                                                               \tag{AP.4}
\]
Its scalar on character $i$ is nonzero precisely when
$\gcd(i,n)=1$. The integral character argument of
A.C2.17 therefore makes $A$ zero on all geometric
two-primary torsion: over the finite unramified
extension of $\Z_2$ containing the $n$-th roots of
unity, primitive Tate summands have rank zero and
each other summand is killed by one factor.
Duality may invert indices but preserves primitivity.

A hypothetical collision of distinct $P,Q\in C(K)$
lifts after finite extension, by A.3, to $P',Q'\in D$
with two-primary difference. Applying $A$ gives the
principal divisor
\[
 \mathcal D=\sum_{S}(-1)^{|S|}
                    (a_SP'-a_SQ'),\qquad
 a_S=\prod_{\ell\in S}a_\ell,                             \tag{AP.5}
\]
where $S$ runs over subsets of the distinct primes
dividing $n$.
The $2^r$ elements $a_S$ are distinct, since their
prime-order factors have pairwise coprime orders.
The two deck orbits are disjoint and free.
Thus (AP.5) has $d$ distinct zeros and $d$ distinct
poles, all simple. Let $f$ be its rational function;
it is separable and has degree $d$.

Over $k=\overline K$ set
\[
              M=k(\gamma f:\gamma\in G),\qquad
              e=[k(D):M],
\]
and let $Z$ be the smooth curve with function field
$M$. Every generator has degree $d$, so (AP.1)
applies. The equality $e=1$ would give
$n+1=g(D)\le(d-1)^2$, contradicting (AP.2).
Hence $e>1$ and is a power of two.

The field $M$ is $G$-stable. The induced action on
$Z$ is faithful: its kernel has order dividing
both $e$ and the odd number $n$. It is free as
well. A nontrivial stabilizer contains an element
of odd prime order $\ell$; its free action on
$D$ partitions the degree-$e$ fibre divisor into
orbits of size $\ell$ with constant multiplicity.
Thus $\ell\mid e$, a contradiction.
It follows that
\[
                 g(Z)-1=n(g(Z/G)-1).
\]
The genus of $Z$ is therefore either one or at
least $n+1$. The bound (AP.1) is less than $n+1$,
by (AP.2), so $g(Z)=1$.

Write $f=h\circ v$ with $v:D\to Z$ of degree $e$.
Choose a zero of the nonconstant function $h$.
Its full inverse image is reduced by (AP.5), so
the zero is simple, the map $v$ is unramified
above it, and the fibre has exactly $e$ geometric
points. Each of $GP'$ and $GQ'$ maps injectively
to $Z$, because an equality of images would
give a stabilizer in its free $G$-action.
The fibre is contained in these two orbits, so
$e\le2$. Thus $e=2$.

The map $D\to Z$ is consequently a separable
quadratic cover of a genus-one curve. Its
involution $\tau$ commutes with $G$, because
$M$ is $G$-stable and the order-two deck group
has no nontrivial automorphism.
Its nonempty fixed set has cardinal divisible
by $n$. Riemann--Hurwitz gives total different
$2n$, and every wild quadratic fixed point
contributes at least two. There are exactly
$n$ fixed points. The ramified
Deuring--Shafarevich formula gives
\[
                         f(D)=2f(Z)+n-1\ge n-1,
\]
contradicting (AP.3)
\citep[Corollary 1.8, pp.~35--36]{A_PrimCrew1984}.
Alternatively, the independent classes
$n[R_i-R_0]$ from the ramification orbit, as
proved in A.C2.17, give the same lower bound
without the formula.

This excludes the collision. Finite-extension
descent and the open-subcurve assertion follow
from A.3 and functoriality to the compactification,
as in A.C2.17; no specialization of the theta
condition is required.

The new point is the support of (AP.5) in only
two deck orbits. Its degree is $2^{\omega(n)}$,
not necessarily four. The controlled field
and reduced zero fibre force the same quadratic
elliptic obstruction for every odd conductor.
This remains a conditional criterion: a single
theta torsion point need not meet all primitive
doubling orbits.
The certificate
\path{code/BC_arbitrary_primitive_prym_checks.py}
and its JSON output record 830 passing checks,
including 500 odd orders and six exact
cyclotomic identities at conductors with at
least three prime divisors.
\end{proof}
